% !TeX root = ../main.tex
\section{Formal Mathematical Model (2-Player Mode)}
\label{app:formal}

This appendix specifies Tactical Formation mathematically,
\emph{independently of any software implementation}, complementing the
presentation through images in Section~\ref{sec:rules} (which is
enough to grasp the rules). The focus is on the core of the game: which
\textbf{direction} a piece moves in and how it \textbf{captures}. The model
rests on three pillars---\textbf{state}, \textbf{action},
\textbf{reward}---and one principle throughout: \emph{the board always has
exactly 28 pieces; capture is a change of side, never removal}. The
extension to the 3--4 player modes is given in Section~\ref{sec:variants}.

\subsection{The board and the 28 pieces}
\label{ssec:board}

The 2-player board consists of 41 cells on two layers:
\begin{align}
\mathcal{B}_{\mathsf V}
  &= \{(i,j)\in\mathbb Z^2 : 0\le i,j\le 4\}
  && \text{(25 vertex cells)}, \label{eq:vert}\\
\mathcal{B}_{\mathsf C}
  &= \{(i{+}\tfrac12,\,j{+}\tfrac12) : 0\le i,j\le 3\}
  && \text{(16 centre cells)}, \label{eq:cent}\\
\mathcal{B}
  &= \mathcal{B}_{\mathsf V}\sqcup\mathcal{B}_{\mathsf C},
  && |\mathcal{B}|=41. \label{eq:board}
\end{align}
Each centre cell lies exactly in the middle of the four surrounding vertex
cells. The type function $\kappa(x)\in\{\mathsf V,\mathsf C\}$ tells whether
cell $x$ is a vertex or a centre. The two sides are $\mathcal C=\{\mathrm
Y,\mathrm G\}$ (Yellow, Green); the opponent of $c$ is written $\bar c$; an
empty cell is denoted $\varnothing$. By convention, Green moves first
(Section~\ref{ssec:turn_order}).

The game has exactly \textbf{28 pieces}; each side starts with 14 pieces
($=10$ vertex pieces $+\,4$ centre pieces). Each piece $p$ carries three
attributes:
\begin{itemize}[noitemsep,topsep=2pt]
  \item \textbf{type} $\kappa(p)\in\{\mathsf V,\mathsf C\}$---\emph{invariant
        for the whole game}: a vertex piece always stays on vertex cells, a
        centre piece on centre cells;
  \item \textbf{position} $x_p$---changes when the piece jumps;
  \item \textbf{side} $c_p\in\{\mathrm Y,\mathrm G\}$---changes when the
        piece is captured.
\end{itemize}

\begin{quote}
\textbf{The 28-piece conservation invariant.} No piece is ever added to or
removed from the board; capture is a \emph{change of side} of an existing
piece. Therefore, at every state
\begin{equation}
N_{\mathrm Y}(s)+N_{\mathrm G}(s)=28,
\label{eq:conservation}
\end{equation}
where $N_c(s)$ is the number of pieces of side $c$.
\end{quote}

We should be clear about the status of this claim: it is an \emph{invariant
that follows directly from the definition}---because capture is defined as
\emph{flipping the colour} of an existing piece (no addition, no removal),
the total piece count obviously does not change---rather than a deep theorem
needing an elaborate proof. Its value lies in its \emph{consequence for
play}: the game never degenerates into a sparse endgame, and all 28 pieces
always stay in play.

The invariant~\eqref{eq:conservation} distinguishes Tactical Formation from
both Checkers (capture $=$ removal) and Reversi (the piece count grows with
each placement), and places it in the same group as Shoulder-Pole Chess and Ming Mang,
two games that also conserve the total number of pieces
(Section~\ref{sec:related}). The starting position (state $0$) consists of two
symmetric blocks of pieces:
\begin{align}
P^{\mathrm Y}_0 &= \{(i,0)\}_{i=0}^{4}\cup\{(i{+}\tfrac12,\tfrac12)\}_{i=0}^{3}\cup\{(i,1)\}_{i=0}^{4},
\label{eq:starty}\\
P^{\mathrm G}_0 &= \{(i,4)\}_{i=0}^{4}\cup\{(i{+}\tfrac12,\tfrac72)\}_{i=0}^{3}\cup\{(i,3)\}_{i=0}^{4}.
\label{eq:startg}
\end{align}

\subsection{State}
\label{ssec:state}

The game state is the pair
\begin{equation}
\boxed{\;s=(b,\mathrm{turn}),\qquad
b:\mathcal B\to\{\mathrm Y,\mathrm G,\varnothing\},\qquad
\mathrm{turn}\in\{\mathrm Y,\mathrm G\}.\;}
\label{eq:state}
\end{equation}
$b$ is the \emph{board configuration} (which cell holds which colour), and
$\mathrm{turn}$ is the side to move. The position and side of all 28
pieces together with the side to move determine the entire state; the game is one of
\emph{perfect information}, with no hidden information. The starting state is
$s_0=(b_0,\mathrm G)$ (by convention Green moves first; in real play the
first mover is agreed among the players) with $b_0$ placed per
\eqref{eq:starty}--\eqref{eq:startg}. The piece-count function $N_c(s)$
satisfies $N_{\mathrm Y}(s_0)=N_{\mathrm G}(s_0)=14$ and the
invariant~\eqref{eq:conservation}.

\subsection{Directions and the jump}
\label{ssec:action}

The set of movement directions depends on the piece type:
\begin{equation}
D_{\mathsf V} = \underbrace{\{(\pm1,0),(0,\pm1)\}}_{\text{4 orthogonal, step }1}
  \cup \underbrace{\{(\pm\tfrac12,\pm\tfrac12)\}}_{\text{4 diagonal, step }\frac12},
\qquad
D_{\mathsf C} = \{(\pm\tfrac12,\pm\tfrac12)\}.
\label{eq:dirs}
\end{equation}
A vertex piece has 8 directions; a centre piece has only 4 diagonal
directions (moving orthogonally by one cell would fall off the centre
layer). A piece at cell $x$ moves along $d\in D_{\kappa(x)}$ by jumping
\emph{over} the adjacent cell $x+d$ and \emph{landing} on $x+2d$. A jump is
valid if and only if
\begin{equation}
\underbrace{x{+}d,\ x{+}2d\in\mathcal B}_{\text{stays on the board}},
\qquad
\underbrace{b(x{+}d)\neq\varnothing}_{\text{springboard occupied (any colour)}},
\qquad
\underbrace{b(x{+}2d)=\varnothing}_{\text{landing cell empty}}.
\label{eq:jump}
\end{equation}
Because $2d$ is an integer vector, the landing cell is automatically of the
same type as $x$ (vertex to vertex, centre to centre). \textbf{The
springboard piece stays put}---it does not move and does not change side.
The set of valid landing cells and the set of moves are
\begin{equation}
\begin{aligned}
\mathrm{Jumps}(x,b)&=\{x{+}2d : d\in D_{\kappa(x)},\ \eqref{eq:jump}\ \text{holds}\},\\
\mathcal A_{\mathrm{move}}(s)&=\{(x,x') : b(x)=\mathrm{turn},\ x'\in\mathrm{Jumps}(x,b)\}.
\end{aligned}
\label{eq:moves}
\end{equation}
The springboard for a vertex piece moving diagonally is a \emph{centre}
cell; the springboard for a centre piece is a \emph{vertex} cell---the two
layers interlock through the diagonals (see the illustration in
Figure~\ref{fig:jumps}).

\subsection{Capturing---the core of the game}
\label{ssec:capture}

In short: from the
landing cell, along each direction, if you see a run of consecutive opponent
pieces immediately followed by one of your own pieces, then the whole run is
\emph{sandwiched} and \emph{changes side} to yours. Capturing uses the same
direction set as movement: a vertex piece considers
$\Delta_{\mathsf V}=D_{\mathsf V}$ (8 directions), a centre piece considers
$\Delta_{\mathsf C}=D_{\mathsf C}$ (4 diagonal directions). The \emph{ray} in
direction $\delta$ from the landing cell $x'$ is the sequence of cells
$x'+\delta,\ x'+2\delta,\ x'+3\delta,\dots$; an orthogonal ray passes only
through vertex cells, while a diagonal ray alternates
centre--vertex--centre, so a diagonal sandwich can capture both vertex and
centre pieces.

Let $b'$ be the board right after the jump (jumped, not yet flipped) and
$c=\mathrm{turn}$. Direction $\delta$ \textbf{captures} when there exists an
integer $m\ge 1$ such that
\begin{equation}
\underbrace{b'(x'+k\delta)=\bar c,\quad k=1,\dots,m}_{m\ \text{consecutive opponent pieces}},
\qquad
\underbrace{b'(x'+(m{+}1)\delta)=c}_{\text{a friendly piece right after the run}}.
\label{eq:clamp}
\end{equation}
Then the flip set of direction $\delta$ is $F_\delta=\{x'+k\delta : 1\le
k\le m\}$. Direction $\delta$ does \emph{not} capture (by convention
$F_\delta=\varnothing$) if the ray---before reaching one of your own
pieces---hits an empty cell, leaves the board, or immediately meets one of
your own pieces (with no opponent piece in between). The set of capturing
directions of move $(x,x')$ is $\Delta(s,x,x')=\{\delta :
F_\delta\neq\varnothing\}$.

\paragraph{The full action.} An action is the triple
\begin{equation}
\boxed{\;a=(x,\ x',\ \delta)\;}
\qquad
\delta\in
\begin{cases}
\Delta(s,x,x') & \text{if a capturing direction exists},\\
\{\varnothing\} & \text{otherwise (no flip)},
\end{cases}
\label{eq:action}
\end{equation}
with $(x,x')\in\mathcal A_{\mathrm{move}}(s)$. \textbf{Rule:} if a sandwich
exists the player \emph{must capture}, and chooses \emph{exactly one
direction}---the other capturing directions do not flip. Every $F_\delta$
can be computed directly from the board-after-the-jump, so the mover always
foresees how many pieces each capturing choice takes.
\paragraph{Example.} \emph{Notation convention in this example:} uppercase
$\mathrm Y,\mathrm G$ denotes \emph{vertex} pieces, lowercase
$\mathrm y,\mathrm g$ denotes \emph{centre} pieces (of the same
Yellow/Green colour); this is merely shorthand to make each piece's layer
visible---the board function $b$ in~\eqref{eq:state} still returns only a
colour, while the layer is the attribute $\kappa(x)$ of the cell. A Green
piece ($c=\mathrm G$) has just landed on $x'=(1,2)$ with the surrounding
configuration: $b'(2,2)=\mathrm Y$, $b'(3,2)=\mathrm G$ on the horizontal
ray; and $b'(1.5,2.5)=\mathrm y$, $b'(2,3)=\mathrm Y$,
$b'(2.5,3.5)=\mathrm y$, $b'(3,4)=\mathrm G$ on the diagonal ray. The two
capturing directions:
\begin{align}
\delta_1=(1,0):\quad & F_{\delta_1}=\{(2,2)\}, && |{\cdot}|=1, \label{eq:ex1}\\
\delta_2=(\tfrac12,\tfrac12):\quad & F_{\delta_2}=\{(1.5,2.5),(2,3),(2.5,3.5)\}, && |{\cdot}|=3. \label{eq:ex2}
\end{align}
The diagonal ray~\eqref{eq:ex2} captures both centre and vertex pieces (2
centre $+$ 1 vertex), illustrating how a diagonal ray alternates between the
two layers. The mover chooses exactly one direction; choosing $\delta_2$,
the three Yellow pieces on the diagonal ray turn Green, while $Y@(2,2)$ on
the horizontal ray \emph{stays Yellow}. This situation is shown visually in
Figure~\ref{fig:capture} and Table~\ref{tab:capture_rays}.

\subsection{Reward}
\label{ssec:reward}

The size of an action is the number of pieces it captures, $|a|:=|F_\delta|$
(equal to $0$ if $\delta=\varnothing$). The per-step reward is zero-sum:
\begin{equation}
\boxed{\;r_c(s,a)=+|a|,\qquad r_{\bar c}(s,a)=-|a|\;}
\qquad (c=\mathrm{turn}).
\label{eq:reward}
\end{equation}
Four direct consequences of \emph{capture $=$
change of side}:
\begin{enumerate}[label=(\arabic*),noitemsep,topsep=2pt]
  \item \textbf{Zero-sum per step:} $r_{\mathrm Y}(s,a)+r_{\mathrm G}(s,a)=0$.
  \item \textbf{Capturing $k$ pieces swings the piece-count difference by
        $2k$:} us $+k$, opponent $-k$ at once.
  \item \textbf{The total is always 28:} the invariant~\eqref{eq:conservation}
        at every state.
  \item \textbf{The scoreboard is the board itself:} along a game
        $s_0,a_0,s_1,\dots,s_t$,
        \begin{equation}
        N_c(s_t)=14+\sum_{k=0}^{t-1} r_c(s_k,a_k);
        \label{eq:scoreboard}
        \end{equation}
        counting each side's pieces at any moment reads off that side's
        cumulative reward.
\end{enumerate}
When the game ends, the winning side receives $+W$ and the losing side $-W$,
where $W>0$ is an abstract constant (each implementation chooses a concrete
value); the two sides' total reward over the whole game is always $0$.

\subsection{The transition function $T$}
\label{ssec:transition}

Given $s=(b,\mathrm{turn})$, $a=(x,x',\delta)$, $c=\mathrm{turn}$. A move
proceeds in three steps.

\noindent\textbf{Step 1---jump.} The piece leaves $x$ and lands on $x'$:
\begin{equation}
b'(y)=\begin{cases}\varnothing & y=x,\\ c & y=x',\\ b(y) & \text{other cells}.\end{cases}
\label{eq:t1}
\end{equation}
\textbf{Step 2---flip along the chosen direction.} Every piece in $F_\delta$
changes to colour $c$ (if $\delta=\varnothing$, nothing flips):
\begin{equation}
b''(y)=\begin{cases}c & y\in F_\delta,\\ b'(y) & \text{other cells}.\end{cases}
\label{eq:t2}
\end{equation}
\textbf{Step 3---pass the turn.}
\begin{equation}
\boxed{\;T(s,a)=(b'',\,\bar c).\;}
\label{eq:t3}
\end{equation}
$T$ is a \emph{deterministic} function---there is no randomness in this game.

\subsection{Terminal conditions}
\label{ssec:terminal}

Side $c$ is \emph{alive} in state $s$ if it still has at least one piece
with a legal jump:
\begin{equation}
\mathrm{alive}(c,s)\iff \exists\,x:\ b(x)=c\ \wedge\ \mathrm{Jumps}(x,b)\neq\varnothing.
\label{eq:alive}
\end{equation}
\begin{equation}
\boxed{\;\text{If the side to move is not alive}
\ \Rightarrow\ \text{that side loses immediately, the other wins.}\;}
\label{eq:terminal}
\end{equation}
The two corresponding ways to lose (Figure~\ref{fig:endgame}): \textbf{out
of pieces} (both the original 14 pieces and any captured pieces have all
changed side) or \textbf{blockaded} (pieces remain but none can jump). A
side that temporarily has no move while it is \emph{not} its turn has not
lost---the condition is checked only at the moment its turn arrives. The
2-player mode has \emph{no draw}.

\subsection{Complexity upper bounds}
\label{ssec:complexity}

The invariance of piece types allows a fairly compact bound
on the size of the state space. Because the piece type is invariant, there
are always exactly 20 vertex pieces (over 25 vertex cells) and 8 centre
pieces (over 16 centre cells); each piece carries one of two colours.
Ignoring reachability, the number of board configurations is bounded by
\begin{equation}
|\mathcal S| \;\le\; 2\cdot
\underbrace{\binom{25}{20}2^{20}}_{\text{vertex layer}}\cdot
\underbrace{\binom{16}{8}\,2^{8}}_{\text{centre layer}}
\;\approx\; 3.7\times10^{17},
\label{eq:statespace}
\end{equation}
(a factor of $2$ for the turn). This is a loose upper bound that includes
states which never arise in real play. It is about three orders of magnitude
below Checkers (about $5\times10^{20}$ positions), which was weakly solved in
2007~\citep{schaeffer2007checkers}. Solving the 2-player mode of Tactical Formation
may therefore be a feasible research direction, although it would still need
substantial computing resources.

As for the \emph{branching factor}, an action is the triple
$(x,x',\delta)$. Because capture is a change of side, one side can hold
almost all 28 pieces, so the bound cannot use the initial 14. The
\emph{kind} of a piece is invariant, however: the board always has exactly
20 vertex pieces and 8 centre pieces. A vertex piece has $\le 8$ landing
cells, and each landing cell offers $\le 8$ sandwich rays of which one is
chosen (a single action if there is none); a centre piece has $\le 4$
landing cells and $\le 4$ rays. Hence the number of actions at a state is
loosely bounded by
\begin{equation}
|\mathcal A(s)| \;\le\; 20\cdot 8\cdot 8 + 8\cdot 4\cdot 4 \;=\; 1408 .
\label{eq:branchbound}
\end{equation}
The jump constraint~\eqref{eq:jump} (requiring both a springboard \emph{and}
an empty landing cell) makes the actual branching factor much smaller; the
measured values are reported in Section~\ref{sec:eval}.

\subsection{The model in summary}
\label{ssec:onelinerule}

Putting it together, the complete model is the tuple
\begin{equation}
\boxed{\;\mathcal G=\big(\mathcal S,\ \mathcal A,\ T,\ r,\ R,\ s_0\big)\;}
\label{eq:game}
\end{equation}
---a \emph{two-player, zero-sum, perfect-information, alternating-turn,
deterministic} game: state~\eqref{eq:state}, action $=$ jump $+$ capturing
direction~\eqref{eq:action}, reward $=$ number of pieces
captured~\eqref{eq:reward}, transition~\eqref{eq:t3},
termination~\eqref{eq:terminal}, starting from $s_0$. Each side maximises
$\sum_t r_c(s_t,a_t)+(\pm W)$---capture as many pieces as possible, and above
all, win.

\paragraph{An agent-based reading.} Besides regarding the two players as the
two decision-making agents, one may regard the 28 pieces as 28 agents split
into two teams: a captured piece moves to the other team, so the two teams
always total 28. Both readings describe the same game $\mathcal G$.
